\documentclass[11pt,a4paper]{article}
\usepackage[T1]{fontenc}
\usepackage[utf8]{inputenc}
\usepackage{amsmath,amssymb,xcolor}
\usepackage[margin=25mm]{geometry}
\usepackage[colorlinks=true,urlcolor=blue]{hyperref}
\definecolor{nexorablue}{HTML}{2459B5}
\setlength{\parindent}{0pt}
\setlength{\parskip}{7pt}
\setlength{\emergencystretch}{2em}
\renewcommand{\familydefault}{\sfdefault}
\begin{document}

{\small\bfseries\color{nexorablue} NEXORA CPGE}\par
{\small MP / NEX-IND-06 / Corrigé}\par
{\LARGE\bfseries Après la disparition du flux extérieur}\par
\vspace{0.25em}\hrule\vspace{0.6em}

\section*{Décroissance}
Avec $\tau=L/R=0{,}050\,\mathrm s$, le courant vaut $i_0e^{-t/\tau}$. Au bout de $\tau$, il vaut environ $0{,}294\,\mathrm A$ et l’énergie devient $(Li_0^2/2)e^{-2}\simeq0{,}0108\,\mathrm J$.

\section*{Dissipation}
L’énergie initiale est $Li_0^2/2=0{,}080\,\mathrm J$ et le courant finit par s’annuler. Toute cette énergie est dissipée par effet Joule.


\end{document}
